
\section{Preliminaries and Problem Formulation}
\label{sec:preliminaries}
 

We formalize the problem of complex reasoning as learning a conditional distribution $p(y|x)$ that maps an input query $x \in \mathcal{X}$ to a final answer $y \in \mathcal{Y}$.
For tasks requiring multi-step logic, learning this direct mapping reliably is difficult, as a single forward pass must implicitly capture all intermediate reasoning. Standard approaches address this by introducing an explicit intermediate rationale that decomposes the reasoning process into tractable steps.

\subsection{Formulation of Reasoning Dynamics}
\label{subsec:formulation}

\paragraph{Explicit Chain-of-Thought.}
In the standard paradigm, the rationale is a sequence of discrete tokens $C = (c_1, c_2, \dots, c_N)$ drawn from a fixed vocabulary $\mathcal{V}$. The joint distribution factorizes as:
\begin{equation}
    p(y, C | x) = p(y | C, x) \prod_{i=1}^N p(c_i | c_{<i}, x).
\end{equation}
Projecting abstract reasoning onto discrete symbols acts as a \textit{lossy compression} of the underlying thought process, fundamentally constrained by the finite vocabulary and the discrete nature of language.% 引用

\paragraph{Latent Chain-of-Thought.}
We investigate \textit{Latent Chain-of-Thought}, where this scaffold is internalized into a trajectory of continuous hidden states $L = (L_1, L_2, \dots, L_T)$.
Each reasoning step $L_t$ consists of a sequence of $g$ continuous vectors, where $g$ denotes the latent granularity:
\begin{equation}
    L_t = (h_{t,1}, \dots, h_{t,g}) \in \mathbb{R}^{g \times d},
\end{equation}
where $d$ is the hidden dimension. The inference marginalizes over the latent trajectory:
\begin{equation}
    p(y|x) = \int p(y|L, x) p(L|x)\,dL.
\end{equation} % latent variable 边缘化公式。



\paragraph{The Representational Interface Gap.}
The contrast between explicit and latent CoT can be viewed as a difference in the bandwidth of the intermediate reasoning interface.
Let \(R\) denote an idealized random variable representing the task-relevant reasoning state underlying a complete solution.
For explicit CoT, this information must pass through \(C \in \mathcal{V}^N\).
Since \(C\) consists of \(N\) vocabulary tokens, its conditional entropy is upper-bounded by
\begin{equation}
    H(C \mid x) \le N \log_2 |\mathcal{V}|.
\end{equation}
In practice, syntactic and semantic constraints of natural language make the empirical entropy of emitted rationales substantially lower.
Equivalently, any information about the underlying reasoning process that is recoverable from the emitted rationale satisfies
\begin{equation}
    I(R; C \mid x) \le H(C \mid x) \le N \log_2 |\mathcal{V}|.
\end{equation}
If the rationale corresponds to \(T\) conceptual reasoning steps with an average of \(\bar{n}\) tokens per step (\(N \approx T \bar{n}\)), this gives
\begin{equation}
    I(R; C \mid x) \lesssim T \bar{n} \log_2 |\mathcal{V}|.
\end{equation}

Latent CoT instead exposes a sequence of continuous intermediate states \(L=(L_1,\ldots,L_T)\), where each \(L_t \in \mathbb{R}^{g \times d}\).
For a simplified upper-bound comparison, we consider a finite-precision or noise-limited discretization \(\tilde{L}=Q(L)\), with at most \(B_{\mathrm{eff}}\) reliable bits per dimension.
Then the information recoverable from the latent interface is bounded by
\begin{equation}
    I(R; \tilde{L} \mid x)
    \le H(\tilde{L} \mid x)
    \le T g d B_{\mathrm{eff}}.
    \label{eq:latent_entropy_bound}
\end{equation}
This suggests that latent reasoning states can encode richer non-linguistic computational features that are inevitably lost under explicit symbolic compression~\citep{zhu2025reasoning,zou2025latentcollaborationmultiagentsystems},  since \(d B_{\mathrm{eff}}\) can substantially exceed \(\log_2 |\mathcal{V}|\).
High-dimensional latent states may carry task-relevant computation, but they may also encode noise, spurious correlations, or representations that drift away from the semantic structure learned during language pre-training. 


\subsection{Supervision Challenges in Latent CoT}
\label{subsec:induction_dilemma}
The theoretical objective of latent Chain-of-Thought (CoT) reasoning is to maximize the mutual information (MI,~\citealt{MI}) between the latent trajectory $L$ and the target answer $y$, such that each latent state $L_t$ captures answer-relevant intermediate information. However, unlike discrete reasoning tokens, latent states have no observable ground truth, and the optimal latent trajectory remains inherently hidden. This makes supervision fundamentally challenging: optimizing solely for the final answer provides no direct guarantee that the learned latent states correspond to meaningful reasoning steps. In high-dimensional continuous spaces, the model may instead rely on shortcut features or spurious representations that are unrelated to valid reasoning~\citep{zhang2025latenttokensthinkcausal}. This setting corresponds to \textbf{outcome supervision (OS)}, where latent states are shaped exclusively by gradients propagated from the final answer.
To mitigate this limitation, existing work increasingly incorporates explicit rationales $C=(c_1,\ldots,c_N)$ as additional supervision signals. We collectively refer to these approaches as \textbf{process supervision (PS)}. Let the rationale be partitioned into $T$ reasoning steps, denoted as $S_1,\ldots,S_T$. We observe that existing process supervision methods can be understood through two complementary dimensions: \textit{trajectory supervision} and \textit{space supervision}.

\paragraph{Trajectory Supervision.}
Trajectory supervision regulates the \emph{temporal dynamics} of latent reasoning.
By requiring the accumulated latent states $L_{\le t}$ to predict future reasoning content, such as the next explicit step $S_{t+1}$ or the future suffix $S_{>t}=(S_{t+1},\ldots,S_T)$, trajectory supervision enforces the sequential predictability of the latent trajectory.
It therefore encourages the continuous flow to follow the logical progression of the explicit chain and provides dense signals for mitigating temporal information decay.

\paragraph{Space Supervision.}
Space supervision regulates the \emph{semantic grounding} of individual latent states.
It encourages each latent state $L_t$ to preserve information about its corresponding explicit step $S_t$.
We distinguish two strategies based on the \textit{space} where the alignment constraint is applied:
1) \textit{Generative Reconstruction (GR):}
This approach performs alignment in the symbolic space.
It employs an auxiliary decoder to recover the corresponding textual reasoning step $S_t$ from $L_t$, thereby encouraging the latent state to retain recoverable semantic content without forcing it into a fixed embedding geometry.
2) \textit{Geometric Compression (GC):}
Conversely, this approach performs alignment directly in the latent space.
Instead of generating raw tokens, it constrains $L_t$ to minimize the geometric distance to an encoded representation of $S_t$, typically using a frozen encoder as a stable semantic anchor.


